%% =============================================================================
%% chapters/06-workflow-energeia.tex
%% PART II: Platform Capabilities and Modules
%% Chapter 6: Workflow & Task Processing Engine (Ponos, Ergon, Praxis, Pragma)
%% =============================================================================

\chapter[Workflow \& Task Processing Engine (Energeia)]{Workflow \& Task Processing Engine\\(\texttt{Energeia})}
\label{chap:workflow_energeia}

The Energeia subsystem is Harmonia's composable, high-performance workflow execution engine. It coordinates discrete clinical processing activities, manages the lifecycle of canonical integration tasks (\texttt{Pragma}), and records end-to-end execution audit trails.

\begin{figure}[htbp]
    \centering
    \begin{adjustbox}{max width=\textwidth}
        %% =============================================================================
%% diagrams/fig-energeia-workflow.tex
%% ArchiMate Energeia Workflow Engine & Camel Sequence View
%% =============================================================================
\begin{tikzpicture}[node distance=1.0cm and 0.8cm, auto]

    % Ponos Ingress
    \node[archi-app-component] (conduit) at (0.0, 7.5) {\archibadge{Conduit}\archiname{Petasos Ingress Conduit}\\\archidesc{Queue Consumer}};
    \node[archi-app-component] (dispatcher) at (5.0, 7.5) {\archibadge{Dispatcher}\archiname{Workflow Dispatcher}\\\archidesc{Praxis Resolver}};
    \node[archi-app-component] (praxis) at (10.5, 7.5) {\archibadge{Orchestrator}\archiname{Praxis Implementation}\\\archidesc{Pipeline Sequence}};

    % Checkpoint Manager & Cache
    \node[archi-app-component] (chk_mgr) at (5.0, 5.2) {\archibadge{Checkpoint Mgr}\archiname{PraxisCheckpointMgr}\\\archidesc{5-Phase Auditor}};
    \node[archi-app-component] (cache_svc) at (10.5, 5.2) {\archibadge{Cache Client}\archiname{PragmaCacheService}\\\archidesc{HotRod Client}};
    \node[archi-data-object] (mneme_box) at (10.5, 3.2) {\archibadge{Data Grid}\archiname{Mneme task-cache}\\\archidesc{Infinispan Store}};

    % Erga Pipeline Steps
    \node[archi-app-process] (erg_step0) at (0.0, 1.0) {\archibadge{Ergon Step 0}\archiname{Adt2FhirMapper}\\\archidesc{Parse HL7 ADT}};
    \node[archi-app-process] (erg_step1) at (3.8, 1.0) {\archibadge{Ergon Step 1}\archiname{ExtractPatient}\\\archidesc{Bundle Unpack}};
    \node[archi-app-process] (erg_step2) at (7.6, 1.0) {\archibadge{Ergon Step 2}\archiname{PatientIdentity}\\\archidesc{EMPI Resolution}};
    \node[archi-app-process] (erg_step3) at (11.4, 1.0) {\archibadge{Ergon Step 3}\archiname{PatientDemographics}\\\archidesc{Attribute Merge}};

    % --- Relationships ---
    % Ingress flow
    \draw[archi-flow] (conduit) -- (dispatcher);
    \draw[archi-flow] (dispatcher) -- (praxis);

    % Praxis to Checkpoint Manager
    \draw[archi-serving] (chk_mgr) -- (praxis);
    \draw[archi-flow] (chk_mgr) -- (cache_svc);
    \draw[archi-flow] (cache_svc) -- (mneme_box);

    % Praxis routes to Erga Steps
    \draw[archi-triggering] (praxis) -- node[left, font=\sffamily\scriptsize]{direct:seq-step-0} (erg_step0);
    \draw[archi-triggering] (erg_step0) -- node[above, font=\sffamily\scriptsize]{step-1} (erg_step1);
    \draw[archi-triggering] (erg_step1) -- node[above, font=\sffamily\scriptsize]{step-2} (erg_step2);
    \draw[archi-triggering] (erg_step2) -- node[above, font=\sffamily\scriptsize]{step-3} (erg_step3);

    % Checkpoint recording flows
    \draw[archi-flow] (erg_step0) -- (chk_mgr);
    \draw[archi-flow] (erg_step1) -- (chk_mgr);
    \draw[archi-flow] (erg_step2) -- (chk_mgr);
    \draw[archi-flow] (erg_step3) -- (chk_mgr);

\end{tikzpicture}

    \end{adjustbox}
    \caption{Energeia Workflow Architecture: Ponos WorkEngine, Praxis Sequences \& Ergon Activities}
    \label{fig:energeia_workflow}
\end{figure}

\section{Ponos Task Processor (\texttt{energeia-ponos})}

\texttt{energeia-ponos} is the runtime worker container responsible for:
\begin{itemize}
    \item Consuming inbound integration tasks from Petasos queues (\texttt{clinical.tasks.inbound}, etc.).
    \item Instantiating and orchestrating Apache Camel execution contexts.
    \item Interacting with the \texttt{Mneme} Infinispan cache grid to store task states and checkpoints.
    \item Enforcing concurrency controls, transaction boundaries, and thread pool scaling.
\end{itemize}

\section{Modular Activity Units (\texttt{energeia-erga})}

The \texttt{energeia-erga} module provides a library of discrete, reusable \texttt{Ergon} activity units. Each Ergon implements a single clinical or technical responsibility adhering to the \texttt{ErgonBase} contract:

\begin{figure}[htbp]
    \centering
    \begin{adjustbox}{max width=\textwidth}
        %% =============================================================================
%% diagrams/fig-ergon-module-architecture.tex
%% ArchiMate 3.2 Ergon Activity Module Architecture & Processing Topology View
%% =============================================================================
\begin{tikzpicture}[node distance=1.0cm and 0.8cm, auto]

    % --- Background Plates (Columns) ---
    \draw[archi-plate-bg] (-0.8, 0.4) rectangle (3.8, 9.4);
    \node[archi-plate-title] at (-0.6, 9.2) {Col 1: Ponos Ingress};

    \draw[archi-plate-bg] (4.2, 0.4) rectangle (8.8, 9.4);
    \node[archi-plate-title] at (4.4, 9.2) {Col 2: Ergon Camel Route};

    \draw[archi-plate-bg] (9.2, 0.4) rectangle (13.8, 9.4);
    \node[archi-plate-title] at (9.4, 9.2) {Col 3: FHIR R5 Processing};

    \draw[archi-plate-bg] (14.2, 0.4) rectangle (18.8, 9.4);
    \node[archi-plate-title] at (14.4, 9.2) {Col 4: Cache \& Handoff};

    % --- Column 1: Ponos Ingress & Conduit Dispatch Layer ---
    \node[archi-app-interface, minimum width=3.8cm] (artemis_queue) at (1.5, 7.8) {%
        \archibadge{App Interface}\archiname{Petasos Queue}\\\archidesc{petasos.queue.inbound.*}%
    };

    \node[archi-app-component, minimum width=3.8cm] (queue_conduit) at (1.5, 4.8) {%
        \archibadge{App Component}\archiname{JMS Conduit}\\\archidesc{QueueToExchangeConduit}%
    };

    \node[archi-app-component, minimum width=3.8cm] (praxis_dispatcher) at (1.5, 1.8) {%
        \archibadge{App Component}\archiname{Praxis Dispatcher}\\\archidesc{TaskProcessorRouteBuilder}%
    };

    % --- Column 2: Ergon Activity Camel Route Lifecycle ---
    \node[archi-app-interface, minimum width=3.8cm] (ergon_input_endpoint) at (6.5, 7.8) {%
        \archibadge{App Interface}\archiname{Input Endpoint}\\\archidesc{direct:activity-<ergonId>}%
    };

    \node[archi-app-component, minimum width=3.8cm] (ergon_route_builder) at (6.5, 4.8) {%
        \archibadge{App Component}\archiname{Ergon Route}\\\archidesc{ErgonBase RouteBuilder DSL}%
    };

    \node[archi-app-process, minimum width=3.8cm] (ergon_process_activity) at (6.5, 1.8) {%
        \archibadge{App Process}\archiname{Activity Hook}\\\archidesc{processActivity(Exchange)}%
    };

    % --- Column 3: Payload Extraction, Transformation & Enrichment ---
    \node[archi-app-process, minimum width=3.8cm] (payload_extractor) at (11.5, 7.8) {%
        \archibadge{App Process}\archiname{Payload Extractor}\\\archidesc{HAPI Terser / Jackson Parser}%
    };

    \node[archi-app-process, minimum width=3.8cm] (fhir_transformer) at (11.5, 4.8) {%
        \archibadge{App Process}\archiname{FHIR R5 Transformer}\\\archidesc{Enrichment \& Security Tags}%
    };

    \node[archi-data-object, minimum width=3.8cm] (pragma_envelope) at (11.5, 1.8) {%
        \archibadge{Data Object}\archiname{Pragma Task Envelope}\\\archidesc{State, Checkpoints, DTO}%
    };

    % --- Column 4: Persistence, Checkpointing & Handoff ---
    \node[archi-app-component, minimum width=3.8cm] (task_cache_service) at (16.5, 7.8) {%
        \archibadge{App Component}\archiname{Task Cache Service}\\\archidesc{HotRod Client / Put}%
    };

    \node[archi-node, minimum width=3.8cm] (mneme_grid) at (16.5, 4.8) {%
        \archibadge{Technology Node}\archiname{Mneme \& Mnemosyne}\\\archidesc{Infinispan Grid \& PostgreSQL}%
    };

    \node[archi-app-interface, minimum width=3.8cm] (ergon_output_endpoint) at (16.5, 1.8) {%
        \archibadge{App Interface}\archiname{Completion Endpoint}\\\archidesc{direct:activity-complete}%
    };

    % --- Flow & Structural Relationships ---
    % Conduit ingestion from Artemis
    \draw[archi-flow]       (artemis_queue) -- node[archi-lbl]{JMS TaskEvent} (queue_conduit);
    \draw[archi-triggering] (queue_conduit) -- (praxis_dispatcher);
    \draw[archi-flow]       (praxis_dispatcher) -- node[archi-lbl]{seq-step-N} (ergon_input_endpoint);

    % Ergon Camel Route execution
    \draw[archi-serving]    (ergon_input_endpoint) -- (ergon_route_builder);
    \draw[archi-triggering] (ergon_route_builder) -- (ergon_process_activity);

    % Processing: Extraction, Transformation, and Envelope State
    \draw[archi-assignment] (ergon_process_activity) -- (payload_extractor);
    \draw[archi-triggering] (payload_extractor) -- (fhir_transformer);
    \draw[archi-access]     (fhir_transformer) -- (pragma_envelope);

    % State Checkpointing & Persistence
    \draw[archi-flow] (fhir_transformer) -- (task_cache_service);
    \draw[archi-flow] (task_cache_service) -- node[archi-lbl]{HotRod} (mneme_grid);

    % Route Completion Handoff
    \draw[archi-triggering] (task_cache_service) -- (ergon_output_endpoint);
    \draw[archi-flow]       (ergon_process_activity) to[bend right=20] (ergon_output_endpoint);

\end{tikzpicture}

    \end{adjustbox}
    \caption{Ergon Modular Activity Architecture and Checkpoint Lifecycle}
    \label{fig:energeia_ergon_architecture}
\end{figure}

\subsection{Standard Ergon Activity Implementations}
\begin{itemize}
    \item \textbf{\texttt{Hl7v2ToFhirPatientErgon}}: Parses raw HL7 v2 ADT segments (PID, PD1, NK1) and constructs canonical FHIR R5 \texttt{Patient} resources.
    \item \textbf{\texttt{PatientIdentityUpdateErgon}}: Queries the Enterprise Master Patient Index (EMPI) in Mneme cache to reconcile Medical Record Numbers (MRNs) into master patient IDs.
    \item \textbf{\texttt{PatientDemographicsUpdateErgon}}: Merges updated demographic data (names, addresses, telecoms) into existing patient records.
    \item \textbf{\texttt{AdtDistributionErgon}}: Evaluates clinical distribution rules and fans out outbound notification tasks to dedicated egress queues.
    \item \textbf{\texttt{AbstractProviderRegistryChangeErgon}}: Base class for provider directory change processors (\texttt{PractitionerChangeErgon}, \texttt{OrganizationChangeErgon}, etc.) performing referential verification.
\end{itemize}

\section{Workflow Orchestration Sequences (\texttt{energeia-praxis})}

\texttt{energeia-praxis} provides the sequence orchestration engine that chains ordered Ergon steps into auditable pipelines.

\begin{figure}[htbp]
    \centering
    \begin{adjustbox}{max width=\textwidth}
        %% =============================================================================
%% diagrams/fig-praxis-workflow-architecture.tex
%% ArchiMate 3.2 Praxis Workflow Sequence Architecture & Pipeline Synthesis View
%% =============================================================================
\begin{tikzpicture}[node distance=1.0cm and 0.8cm, auto]

    % --- Column 1: Declarative Metadata & Seeding Tier ---
    \node[archi-data-object] (praxis_def_json) at (0, 7.5) {%
        \archibadge{Data Object}\archiname{Praxis Definition}\\\archidesc{JSON Descriptor / Config}%
    };

    \node[archi-app-component] (praxis_service) at (0, 4.5) {%
        \archibadge{App Component}\archiname{Praxis Service}\\\archidesc{PraxisService HotRod Client}%
    };

    \node[archi-node] (sequence_cache) at (0, 1.5) {%
        \archibadge{Technology Node}\archiname{Task Sequence Cache}\\\archidesc{tasksequence-cache Grid}%
    };

    % --- Column 2: Dynamic Discovery & CDI Injection Tier ---
    \node[archi-app-component] (sequence_loader) at (4.5, 7.5) {%
        \archibadge{App Component}\archiname{Sequence Loader}\\\archidesc{TaskSequenceLoader CDI}%
    };

    \node[archi-app-component] (cdi_ergon_beans) at (4.5, 4.5) {%
        \archibadge{App Component}\archiname{CDI Ergon Resolution}\\\archidesc{Instance<ErgonBase> Injection}%
    };

    \node[archi-app-component] (praxis_impl) at (4.5, 1.5) {%
        \archibadge{App Component}\archiname{Praxis Implementation}\\\archidesc{Praxis / Route Builder}%
    };

    % --- Column 3: Dynamic Camel Pipeline Topology Tier ---
    \node[archi-app-interface] (pipeline_start) at (9.0, 7.5) {%
        \archibadge{App Interface}\archiname{Pipeline Ingress}\\\archidesc{direct:seq-<praxisId>-start}%
    };

    \node[archi-app-process] (pipeline_steps) at (9.0, 4.5) {%
        \archibadge{App Process}\archiname{Chained Erga Steps}\\\archidesc{direct:seq-step-0 \dots\ step-N}%
    };

    \node[archi-app-interface] (pipeline_complete) at (9.0, 1.5) {%
        \archibadge{App Interface}\archiname{Pipeline Egress}\\\archidesc{direct:seq-<praxisId>-complete}%
    };

    % --- Column 4: Checkpointing, Storage & Runtime Operations Tier ---
    \node[archi-app-component] (checkpoint_manager) at (13.5, 7.5) {%
        \archibadge{App Component}\archiname{Checkpoint Manager}\\\archidesc{PraxisCheckpointManager}%
    };

    \node[archi-node] (pragma_cache) at (13.5, 4.5) {%
        \archibadge{Technology Node}\archiname{Pragma Audit Cache}\\\archidesc{pragma-cache Infinispan}%
    };

    \node[archi-app-component] (mgmt_resource) at (13.5, 1.5) {%
        \archibadge{App Component}\archiname{Workflow Ops REST}\\\archidesc{/workflow \& ponos-cli}%
    };

    % --- Flow & Structural Relationships ---
    % Declarative loading and cache
    \draw[archi-access] (praxis_service) -- (praxis_def_json);
    \draw[archi-flow] (praxis_service) -- node[right, font=\tiny] {HotRod CRUD} (sequence_cache);
    \draw[archi-serving] (praxis_service) -- (sequence_loader);

    % Dynamic resolution and instantiation
    \draw[archi-triggering] (sequence_loader) -- (cdi_ergon_beans);
    \draw[archi-assignment] (cdi_ergon_beans) -- (praxis_impl);
    \draw[archi-triggering] (sequence_loader) -- (praxis_impl);

    % Dynamic Camel Route Generation
    \draw[archi-triggering] (praxis_impl) -- node[above, font=\tiny] {createSequenceRoute} (pipeline_start);
    \draw[archi-flow] (pipeline_start) -- (pipeline_steps);
    \draw[archi-flow] (pipeline_steps) -- (pipeline_complete);

    % Checkpointing and State Persistence
    \draw[archi-flow] (pipeline_steps) -- node[above, font=\tiny] {Audit Events} (checkpoint_manager);
    \draw[archi-flow] (checkpoint_manager) -- node[right, font=\tiny] {Save Checkpoint} (pragma_cache);

    % Operations & Dynamic Reload
    \draw[archi-serving] (mgmt_resource) -- (sequence_loader);
    \draw[archi-access] (mgmt_resource) -- (sequence_cache);

\end{tikzpicture}

    \end{adjustbox}
    \caption{Praxis Dynamic Route Synthesis and Multi-Stage Pipeline Execution}
    \label{fig:energeia_praxis_architecture}
\end{figure}

\subsection{Dynamic Route Synthesis \& Checkpointing}
Praxis constructs Camel routes dynamically using CDI bean discovery. At each Ergon boundary, Praxis creates an immutable \texttt{PragmaCheckpoint} capturing:
\begin{equation}
\begin{aligned}
\langle\text{Checkpoint}\rangle = (\text{checkpointId}, \text{pragmaId}, \text{ergonId}, \text{stageName}, \\
                                  &\quad\text{timestamp}, \text{status}, \text{statusMessage}, \text{stepIndex})
\end{aligned}
\end{equation}
Checkpoints are written to the \texttt{Mneme} cache grid, ensuring complete visibility and audit lineage for in-flight tasks.

\newpage
